\chapter{Python Implementation and Reproducible Terminal Results}
\label{chap:implementation}
This chapter connects the preceding mathematical derivations to the actual Python project provided with this report. The listings below are \emph{excerpts of the project source files}, not pseudocode. The console blocks were captured by executing the supplied scripts on 24 October 2026. These are software simulations with a controllable fault observer, not physical STM32 measurements.

\section{Repository structure and execution}
The principal modules are \texttt{aes.py} (AES operations), \texttt{sm4.py} (SM4 operations), \texttt{piret.py} (DFA and round peeling), \texttt{giraud.py} (single-bit candidate filtering), and \texttt{run.py} (combined demonstration). The repository stores these modules in \texttt{src/}; \texttt{src/gui.py} provides the interactive interface. The commands below are run from the repository root.
\begin{lightcode}
python src/giraud.py --algorithm aes --known-bit
python src/giraud.py --algorithm sm4 --known-bit
python src/piret.py --algorithm sm4
\end{lightcode}
The project also supports an unknown-bit model by omitting \texttt{--known-bit}. The \texttt{--known-bit} option is used for the reproducible outputs reproduced here.

\section{AES-128 implementation: S-box, diffusion, and key expansion}
The following code implements AES finite-field multiplication and column mixing. This is the operation responsible for the characteristic four-byte differential exploited by Piret--Quisquater.
\par\smallskip\noindent\textbf{AES finite-field arithmetic, from aes.py}\par
\begin{lightcode}
def xtime(a):
    """Multiply a byte by x (02) in GF(2^8), reducing by polynomial 0x11B."""
    result = a << 1

    # If the x^8 term overflows, reduction contributes 0x1B to the low byte.
    if a & 0x80:
        result ^= 0x1B

    return result & 0xFF


def gf_mul(a, b):
    """Multiply bytes in the AES field; field addition is XOR, not integer +."""
    result = 0

    for _ in range(8):
        if b & 1:
            result ^= a

        # Shift-and-add multiplication, using xtime for polynomial reduction.
        a = xtime(a)
        b >>= 1

    return result


# ============================================================================
# 5. ROUND TRANSFORMATIONS: ALL MATRIX OPERATIONS ARE IN PLACE
# ============================================================================
\end{lightcode}
\par\smallskip\noindent\textbf{AES MixColumns, from aes.py}\par
\begin{lightcode}
def mix_column(column):
    """Return M * column in GF(2^8), without modifying the four input bytes.

    Copying a column before mixing matters: every output byte must use the same
    original input, never a value already overwritten earlier in the operation.
    """
    result = []
    for coefficients in MIX_MATRIX:
        value = 0
        for coefficient, byte in zip(coefficients, column):
            value ^= gf_mul(coefficient, byte)
        result.append(value)
    return result


def mix_columns(state):
    """Mix each column in place; the normal algorithm never prints output."""
    for column in range(4):
        mixed = mix_column([state[row][column] for row in range(4)])
        for row in range(4):
            state[row][column] = mixed[row]
\end{lightcode}
\par\smallskip\noindent\textbf{AES round-key expansion, from aes.py}\par
\begin{lightcode}
def key_expansion(key):
    """Expand the 16-byte master key into eleven independent round-key matrices.

    Each new word XORs the word four positions back with the previous word.
    Every fourth word first applies RotWord, SubWord and the round constant.
    """
    validate_block(key, "AES-128 key")
    # Four initial words become 44 words: K0 plus ten encryption round keys.
    words = [list(key[4 * i:4 * i + 4]) for i in range(4)]

    for i in range(4, 4 * (N_ROUNDS + 1)):
        temp = words[i - 1].copy()

        if i % 4 == 0:
            # Rcon changes only the first byte of the transformed key word.
            temp = sub_word(rot_word(temp))
            temp[0] ^= RCON[i // 4]

        new_word = [words[i - 4][j] ^ temp[j] for j in range(4)]

        words.append(new_word)

    round_keys = []

    for round_number in range(N_ROUNDS + 1):
        # Convert four consecutive words into the same layout as the data state.
        key_bytes = []

        for i in range(4):
            key_bytes.extend(words[round_number * 4 + i])

        round_keys.append(bytes_to_state(key_bytes))

    return round_keys


# ============================================================================
# 7. NORMAL ENCRYPTION / DECRYPTION AND OPTIONAL EXTERNAL OBSERVATION
# ============================================================================
\end{lightcode}

\section{SM4 implementation: rotations, S-boxes, and 32 rounds}
The SM4 code keeps the encryption diffusion $L$ distinct from the key-schedule diffusion $L'$. The observer hook allows the test harness to modify an intermediate value at a chosen round.
\par\smallskip\noindent\textbf{SM4 rotations, diffusion, and substitution, from sm4.py}\par
\begin{lightcode}
def rotl(val: int, count: int) -> int:
    """Rotate a 32-bit unsigned word left; overflow wraps into the low bits."""
    return ((val << count) & WORD_MASK) | (val >> (32 - count))


def linear(b: int) -> int:
    """Data diffusion L: XOR the word with rotations by 2, 10, 18 and 24.

    L is XOR-linear: L(a XOR b) = L(a) XOR L(b). The SM4 DFA uses this
    property to invert an observed output difference before testing the S-box.
    """
    return b ^ rotl(b, 2) ^ rotl(b, 10) ^ rotl(b, 18) ^ rotl(b, 24)


def linear_key(b: int) -> int:
    """Key-schedule diffusion L-prime: rotations 13 and 23, not data-map L."""
    return b ^ rotl(b, 13) ^ rotl(b, 23)


def sbox(a: int) -> int:
    """Apply tau: substitute all four bytes, preserving big-endian word order."""
    input_bytes = struct.pack('>I', a)
    substituted_bytes = bytes(SBOX[byte] for byte in input_bytes)
    return struct.unpack('>I', substituted_bytes)[0]


# ============================================================================
# 4. SINGLE-ROUND RECURRENCES
# ============================================================================
\end{lightcode}
\par\smallskip\noindent\textbf{SM4 round function and round-key generation, from sm4.py}\par
\begin{lightcode}
def round(x: list[int], rk: int, observer=None, round_index=0):
    """Return the next rolling window without modifying the supplied words.

    A = X[i+1] XOR X[i+2] XOR X[i+3] XOR rk[i]
    X[i+4] = X[i] XOR L(tau(A))
    """
    sbox_input = observe(observer, round_index, "Sbox input", x[1] ^ x[2] ^ x[3] ^ rk)
    substituted = observe(observer, round_index, "Sbox output", sbox(sbox_input))
    diffused = observe(observer, round_index, "Linear output", linear(substituted))
    new_word = observe(observer, round_index, "New word", x[0] ^ diffused)
    return x[1], x[2], x[3], new_word


def generate_round_key(k: list[int], round: int) -> int:
    """Compute K[i+4] from four consecutive schedule words and CK[i].

    The parameter name round is retained for callers using keyword arguments.
    Encryption and expansion share tau, but expansion must use linear_key.
    """
    return k[0] ^ linear_key(sbox(k[1] ^ k[2] ^ k[3] ^ CK[round]))


# ============================================================================
# 5. KEY EXPANSION: FOUR MASTER WORDS TO 32 ROUND KEYS
# ============================================================================


def key_expansion(key: bytes) -> list[int]:
    """Return rk0..rk31; each is one new word from the key recurrence."""
    if len(key) != BLOCK_SIZE:
        raise ValueError("SM4 key must contain exactly 16 bytes.")
    master_words = struct.unpack(">4I", key)
    schedule_words = [master_words[i] ^ FK[i] for i in range(4)]

    for i in range(N_ROUNDS):
        # Keep previous words for readability; the recurrence consumes a window.
        next_key = generate_round_key(schedule_words[i:i+4], i)
        schedule_words.append(next_key)

    # The initial four whitened words are not encryption round keys.
    return schedule_words[4:]


# ============================================================================
# 6. SHARED BLOCK PROCESSING AND PUBLIC ENCRYPT/DECRYPT WRAPPERS
# ============================================================================
\end{lightcode}

\section{AES Giraud: filtering last-round key candidates}
For each possible final-round key byte, the implementation undoes the final AddRoundKey and S-box, and tests whether the recovered input difference is a single bit. The solver does not receive the master key.
\par\smallskip\noindent\textbf{AES Giraud candidate test, from giraud.py}\par
\begin{lightcode}
def aes_candidates(correct, faulty, known_mask=None):
    """Return the affected output position and its compatible K10 byte values.

    For each guess k:
      delta = InvSbox[C[j] XOR k] XOR InvSbox[F[j] XOR k]
    Keep k when delta has exactly one set bit, or equals the known bit mask.
    """
    aes.validate_block(correct, "Correct ciphertext")
    aes.validate_block(faulty, "Faulty ciphertext")
    if known_mask is not None:
        validate_mask(known_mask)

    # Without final-round MixColumns, only one ciphertext byte should change.
    changed = [index for index in range(16) if correct[index] != faulty[index]]
    if len(changed) != 1:
        raise ValueError("AES Giraud requires exactly one changed ciphertext byte.")
    position = changed[0]
    allowed = SINGLE_BITS if known_mask is None else (known_mask,)

    candidates = set()
    for guess in range(256):
        clean_input = AES_INVERSE_SBOX[correct[position] ^ guess]
        faulty_input = AES_INVERSE_SBOX[faulty[position] ^ guess]
        difference = clean_input ^ faulty_input
        if difference in allowed:
            candidates.add(guess)
    return position, candidates
\end{lightcode}
\subsection{Observed AES terminal output}
The following excerpt is from an actual execution of \texttt{python src/giraud.py --algorithm aes --known-bit}. It shows the filtering process and final verification.
\par\smallskip\noindent\textbf{Recorded AES Giraud execution excerpt}\par
\begin{lightcode}

============================================================================
GIRAUD / AES-128
Correct ciphertext: 69C4E0D86A7B0430D8CDB78070B4C55A
Bit known to solver: True
AES: require HW(InvSbox[C XOR k] XOR InvSbox[F XOR k]) = 1.

Key byte 0, observation 1, simulator mask 01
Correct: 69 C4 E0 D8 6A 7B 04 30 D8 CD B7 80 70 B4 C5 5A
Faulty:  76 C4 E0 D8 6A 7B 04 30 D8 CD B7 80 70 B4 C5 5A
XOR:     1F 00 00 00 00 00 00 00 00 00 00 00 00 00 00 00
Candidate count: 256 -> 4
Survivors: 0A 0C 13 15

Key byte 0, observation 2, simulator mask 02
Correct: 69 C4 E0 D8 6A 7B 04 30 D8 CD B7 80 70 B4 C5 5A
Faulty:  1B C4 E0 D8 6A 7B 04 30 D8 CD B7 80 70 B4 C5 5A
XOR:     72 00 00 00 00 00 00 00 00 00 00 00 00 00 00 00
Candidate count: 4 -> 1
Survivors: 13

Key byte 13, observation 1, simulator mask 01
Correct: 69 C4 E0 D8 6A 7B 04 30 D8 CD B7 80 70 B4 C5 5A
Faulty:  69 C4 E0 D8 6A 7B 04 30 D8 CD B7 80 70 83 C5 5A
XOR:     00 00 00 00 00 00 00 00 00 00 00 00 00 37 00 00
Candidate count: 256 -> 2
Survivors: 1C 2B

Key byte 13, observation 2, simulator mask 02
Correct: 69 C4 E0 D8 6A 7B 04 30 D8 CD B7 80 70 B4 C5 5A
Faulty:  69 C4 E0 D8 6A 7B 04 30 D8 CD B7 80 70 7B C5 5A
XOR:     00 00 00 00 00 00 00 00 00 00 00 00 00 CF 00 00

... [intermediate observations omitted from this printed excerpt]

XOR:     00 00 00 9A 00 00 00 00 00 00 00 00 00 00 00 00
Candidate count: 2 -> 1
Survivors: 7F

Recovered K10: 13111D7FE3944A17F307A78B4D2B30C5
Recovered master key: 000102030405060708090A0B0C0D0E0F
Verification: PASS / 32 faulty outputs
\end{lightcode}

\section{SM4 Giraud-inspired filtering}
SM4 requires an extra step: invert the linear transformation $L$ to isolate the S-box output difference. The following code tests candidate S-box inputs, then translates them into candidates for a byte of the last-round subkey using the three unchanged state words.
\par\smallskip\noindent\textbf{SM4 S-box differential candidate filter, from giraud.py}\par
\begin{lightcode}
def sm4_candidates(correct, faulty, known_mask=None):
    """Return affected byte, round-key guesses, input guesses and DeltaB.

    Inputs may be actual final ciphertexts or peeled views at an earlier round.

    1. Observe DeltaX35 from the first output word.
    2. Undo linear diffusion: DeltaB = L-inverse(DeltaX35).
    3. Keep a if S(a) XOR S(a XOR e) matches the observed byte difference.
    4. Convert a into a key guess using the three unchanged output words.
    """
    aes.validate_block(correct, "Correct ciphertext")
    aes.validate_block(faulty, "Faulty ciphertext")
    if known_mask is not None:
        validate_mask(known_mask)

    # SM4 outputs X35 || X34 || X33 || X32; only X35 may change in this model.
    if correct[4:] != faulty[4:]:
        raise ValueError("SM4 last-round model requires three output words unchanged.")
    delta_x35 = int.from_bytes(correct[:4], "big") ^ int.from_bytes(faulty[:4], "big")
    delta_b = inverse_linear(delta_x35)
    differences = delta_b.to_bytes(4, "big")
    changed = [index for index in range(4) if differences[index] != 0]
    if len(changed) != 1:
        raise ValueError("SM4 inverse-L difference must affect one S-box byte.")
    position = changed[0]
    allowed = SINGLE_BITS if known_mask is None else (known_mask,)

    # The input value a is unknown. Test every value against each permitted bit.
    inputs = set()
    for guess in range(256):
        for error in allowed:
            difference = sm4.SBOX[guess] ^ sm4.SBOX[guess ^ error]
            if difference == differences[position]:
                inputs.add(guess)
                break

    # A = X32 XOR X33 XOR X34 XOR rk31; all three data words are observable.
    x34 = int.from_bytes(correct[4:8], "big")
    x33 = int.from_bytes(correct[8:12], "big")
    x32 = int.from_bytes(correct[12:16], "big")
    known_byte = ((x32 ^ x33 ^ x34) >> (24 - 8 * position)) & 0xFF
    candidates = {guess ^ known_byte for guess in inputs}
    return position, candidates, inputs, delta_b
\end{lightcode}
\subsection{Observed SM4 terminal output}
The first fault produces a 32-bit ciphertext difference; the inverse-linear operation reduces it to a one-byte S-box difference. A second independent observation reduces the key-byte candidates from two to one.
\par\smallskip\noindent\textbf{Recorded SM4 Giraud-inspired execution excerpt}\par
\begin{lightcode}

============================================================================
GIRAUD-INSPIRED / SM4
Correct ciphertext: 681EDF34D206965E86B3E94F536E4246
Bit known to solver: True
SM4: invert L, then test S(a) XOR S(a XOR e) for one-bit errors.

Round 31, key byte 0, observation 1
Correct: 68 1E DF 34 D2 06 96 5E 86 B3 E9 4F 53 6E 42 46
Faulty:  1B 50 E2 09 D2 06 96 5E 86 B3 E9 4F 53 6E 42 46
XOR:     73 4E 3D 3D 00 00 00 00 00 00 00 00 00 00 00 00
Undo recovered later rounds: {}
Correct: 68 1E DF 34 D2 06 96 5E 86 B3 E9 4F 53 6E 42 46
Faulty:  1B 50 E2 09 D2 06 96 5E 86 B3 E9 4F 53 6E 42 46
XOR:     73 4E 3D 3D 00 00 00 00 00 00 00 00 00 00 00 00
Candidates: 256 -> 2

Round 31, key byte 0, observation 2
Correct: 68 1E DF 34 D2 06 96 5E 86 B3 E9 4F 53 6E 42 46
Faulty:  ED 9D D9 32 D2 06 96 5E 86 B3 E9 4F 53 6E 42 46
XOR:     85 83 06 06 00 00 00 00 00 00 00 00 00 00 00 00
Undo recovered later rounds: {}
Correct: 68 1E DF 34 D2 06 96 5E 86 B3 E9 4F 53 6E 42 46
Faulty:  ED 9D D9 32 D2 06 96 5E 86 B3 E9 4F 53 6E 42 46
XOR:     85 83 06 06 00 00 00 00 00 00 00 00 00 00 00 00
Candidates: 2 -> 1

Round 31, key byte 1, observation 1
Correct: 68 1E DF 34 D2 06 96 5E 86 B3 E9 4F 53 6E 42 46
Faulty:  AD AB AF F1 D2 06 96 5E 86 B3 E9 4F 53 6E 42 46
XOR:     C5 B5 70 C5 00 00 00 00 00 00 00 00 00 00 00 00
Undo recovered later rounds: {}

... [intermediate observations omitted from this printed excerpt]

Candidates: 2 -> 1
Recovered rk31: 9124A012
Recovered rk30: 01CF72E5
Recovered rk29: 62293496
Recovered rk28: 428D3654
Recovered master key: 0123456789ABCDEFFEDCBA9876543210
Verification: PASS / 32 faulty outputs / no reference subkeys
\end{lightcode}

\section{Recovering SM4 rounds 31 through 28}
The source project implements an additional \emph{multi-round simulated attack}, beyond the last-round example described earlier in this report. Once $rk_{31}$ is recovered, its value is used to undo the last round in both correct and faulty ciphertexts. A separate controlled fault at round 30 can then be analyzed as a last-round differential on the reconstructed state. The same process is repeated for rounds 29 and 28.

Crucially, the algorithm does \emph{not} supply reference round keys to the recovery function. It uses previously recovered keys to peel later rounds. However, its fault oracle must be capable of injecting faults at each of the four targeted rounds, a strong capability not established by these software tests on physical hardware.
\par\smallskip\noindent\textbf{Undoing already recovered SM4 rounds, from piret.py}\par
\begin{lightcode}
def peel_sm4(ciphertext, target_round, recovered_keys):
    """Undo only rounds whose subkeys have ALREADY been recovered by DFA.

    Return a reversed four-word output at the end of target_round. This lets
    the same last-round equations analyze rounds 30, 29 and 28 in succession.
    No reference round keys or secret intermediate values are accepted.
    """
    aes.validate_block(ciphertext, "SM4 ciphertext")
    if not 28 <= target_round <= 31:
        raise ValueError("Full SM4 recovery targets rounds 28..31.")
    words = list(struct.unpack(">4I", ciphertext))[::-1]
    for number in range(31, target_round, -1):
        if number not in recovered_keys:
            raise ValueError(f"Cannot undo round {number} before recovering its subkey.")
        previous = words[3] ^ sm4.linear(sm4.sbox(words[0] ^ words[1] ^ words[2] ^ recovered_keys[number]))
        words = [previous, words[0], words[1], words[2]]
    return struct.pack(">4I", *words[::-1])
\end{lightcode}
\par\smallskip\noindent\textbf{Sequential recovery of four SM4 subkeys, from piret.py}\par
\begin{lightcode}
def recover_sm4_master_key(correct, fault_oracle, masks=(1, 2, 3, 4, 5, 6, 7, 8), start_byte=0, candidate_filter=None, observer=None):
    """Recover rk31 -> rk28 and reverse the schedule, using outputs only.

    fault_oracle(round, byte, mask) returns a faulty ciphertext. It represents
    a target capable of a controlled fault in EACH of the final four rounds.
    The solver receives no master key, schedule, or reference subkeys.
    The default filter uses known masks; Giraud supplies its single-bit filter.
    """
    aes.validate_block(correct, "Correct SM4 ciphertext")
    masks = tuple(masks)
    if not masks or len(set(masks)) != len(masks) or any(not 1 <= mask <= 255 for mask in masks) or not 0 <= start_byte < 4:
        raise ValueError("Use distinct nonzero byte masks and start_byte 0..3.")
    keys, history, counts, used = {}, [], {}, 0
    for number in range(31, 27, -1):
        clean_view = peel_sm4(correct, number, keys)
        subkey = bytearray(4)
        for byte in [(start_byte + offset) % 4 for offset in range(4)]:
            survivors = None
            for observation, mask in enumerate(masks, 1):
                faulty = bytes(fault_oracle(number, byte, mask))
                faulty_view = peel_sm4(faulty, number, keys)
                if candidate_filter is None:
                    fresh, inputs, delta = sm4_byte_candidates(clean_view, faulty_view, byte, mask)
                else:
                    position, fresh, inputs, delta = candidate_filter(clean_view, faulty_view, mask)
                    if position != byte:
                        raise ValueError("Fault observation does not match the requested byte.")
                before = 256 if survivors is None else len(survivors)
                survivors = set(fresh) if survivors is None else survivors & fresh
                if not survivors:
                    raise ValueError("SM4 candidate intersection is empty; fault model is incompatible.")
                used += 1
                counts[f"{number}:{byte}"] = len(survivors)
                event = dict(round=number, byte=byte, observation=observation, mask=mask, before=before, pair_candidates=len(fresh), remaining=len(survivors), survivors=sorted(survivors), inputs=sorted(inputs), delta_b=delta, correct=correct.hex(), faulty=faulty.hex(), clean_view=clean_view.hex(), faulty_view=faulty_view.hex(), peeled_with={str(i): f"{value:08X}" for i, value in keys.items()})
                history.append(event)
                if observer is not None:
                    observer(dict(event))
                if len(survivors) == 1:
                    subkey[byte] = next(iter(survivors))
                    break
            else:
                return dict(master_key=None, round_keys=keys, history=history, candidate_counts=counts, used=used, complete=False, blocked_round=number)
        keys[number] = int.from_bytes(subkey, "big")
    master = reverse_sm4([keys[number] for number in range(28, 32)])
    return dict(master_key=master, round_keys=keys, history=history, candidate_counts=counts, used=used, complete=True)
\end{lightcode}
\section{Reversing the SM4 key schedule}
After recovering four consecutive 32-bit subkeys, the code reverses the key schedule and removes the initial FK constants.
\par\smallskip\noindent\textbf{Reverse SM4 key expansion, from piret.py}\par
\begin{lightcode}
def reverse_sm4(final_keys, observer=None):
    """Four consecutive keys rk28..rk31 are REQUIRED, not merely rk31."""
    if len(final_keys) != 4:
        raise ValueError("SM4 reversal requires four consecutive final subkeys.")
    words = [None] * 32 + list(final_keys)
    for index in range(31, -1, -1):
        value = words[index+1] ^ words[index+2] ^ words[index+3] ^ sm4.CK[index]
        words[index] = words[index+4] ^ sm4.linear_key(sm4.sbox(value))
        if observer:
            observer(index, words[index])
    return struct.pack(">4I", *(words[i] ^ sm4.FK[i] for i in range(4)))
\end{lightcode}

\subsection{Independent terminal validation}
The controlled-byte model also runs independently in \texttt{piret.py}. The following excerpt shows the real output of the supplied code:
\par\smallskip\noindent\textbf{Recorded SM4 controlled-byte DFA execution excerpt}\par
\begin{lightcode}

============================================================================
SM4: NORMAL ENCRYPTION
Round 00: A=F002C39E B=18E992B1 L(B)=26D99622 new=27FAD345
Round 01: A=EE1470AC B=C62AD40D L(B)=2820815D new=A18B4CB2
Round 02: A=AA4F1C7D B=23A849C4 L(B)=EF1D58B2 new=11C1E22A
Round 03: A=EC195DAA B=C5445623 L(B)=BA47D0FE new=CC13E2EE
Round 04: A=4A2A2C82 B=3C0BED8A L(B)=DF868890 new=F87C5BD5
Round 05: A=52C45770 B=81BB8BD4 L(B)=92A94BE5 new=33220757
Round 06: A=B1F637DF B=DB4D95B0 L(B)=663520BD new=77F4C297
Round 07: A=98DCAF44 B=ADB86FF3 L(B)=B6851005 new=7A96F2EB
Round 08: A=9B600757 B=301EB78B L(B)=DFA69BAA new=27DAC07F
Round 09: A=9DE0BDBE B=8C896C5B L(B)=71FF084E new=42DD0F19
Round 10: A=DC966E60 B=B81A781E L(B)=CF511895 new=B8A5DA02
Round 11: A=A3474E33 B=2EBA4FA9 L(B)=EAE7D511 new=907127FA
Round 12: A=0381926D B=FEBF5D21 L(B)=AC4FEBFC new=8B952B83
Round 13: A=939943CC B=A493FC1F L(B)=96F67340 new=D42B7C59
Round 14: A=8B75648F B=CED363A1 L(B)=97598233 new=2FFC5831
Round 15: A=60069A4A B=1E3D323C L(B)=66EF4F72 new=F69E6888
Round 16: A=DC69F8C8 B=B8227911 L(B)=24B11947 new=AF2432C4
Round 17: A=05F35DDE B=E1EC56B4 L(B)=3935B407 new=ED1EC85E
Round 18: A=7823DB74 B=4CF4129F L(B)=7A5FE213 new=55A3BA22
Round 19: A=85BD0481 B=C76CCCBF L(B)=E4D57022 new=124B18AA
Round 20: A=42680EC9 B=A7252CD9 L(B)=C5C3409B new=6AE7725F
Round 21: A=B5C5D18D B=FDDDC161 L(B)=19D569A7 new=F4CBA1F9
Round 22: A=4B725B6C B=19464B01 L(B)=486E4032 new=1DCDFA10
Round 23: A=1A00D498 B=13D6A5AD L(B)=3DBD1EA9 new=2FF60603
Round 24: A=716B85E6 B=003BC777 L(B)=85153D83 new=EFF24FDC
Round 25: A=C0E8A67F B=8D65CA9E L(B)=9B2FCA8C new=6FE46B75
Round 26: A=A1C2A841 B=F6AFC007 L(B)=94F9AABD new=893450AD
Round 27: A=F85A7885 B=790F4CC7 L(B)=5465894F new=7B938F4C
Round 28: A=DFCE82C0 B=B05A8A8D L(B)=BC9C0D9A new=536E4246

... [intermediate observations omitted from this printed excerpt]

Survivors: 54
DFA RECOVERED rk31 = 9124A012 / PASS
DFA RECOVERED rk30 = 01CF72E5 / PASS
DFA RECOVERED rk29 = 62293496 / PASS
DFA RECOVERED rk28 = 428D3654 / PASS
SM4 MASTER KEY: 0123456789ABCDEFFEDCBA9876543210 / PASS
32 faulty outputs. All four subkeys came from DFA, not references.
\end{lightcode}

\section{Interpretation of the results}
The terminal output demonstrates successful \emph{software-model} recovery of the AES-128 master key and the SM4 master key. In the SM4 Giraud-inspired known-bit demonstration, the program reports 32 faulty outputs, the four recovered round keys, and a final successful verification. The corresponding values are:
\begin{align*}
rk_{31}&=\hex{9124A012}, & rk_{30}&=\hex{01CF72E5},\\
rk_{29}&=\hex{62293496}, & rk_{28}&=\hex{428D3654},\\
MK&=\hex{01234567\allowbreak{}89ABCDEF\allowbreak{}FEDCBA98\allowbreak{}76543210}.
\end{align*}
This extends the earlier pedagogical last-round discussion: the supplied project \emph{does} demonstrate four-subkey recovery under its simulated, selectable-round fault oracle. It does \emph{not} establish that the same fault precision or success rate is achievable through ChipWhisperer glitches on the STM32 target.
